% TEMPLATE-MILC-v3.tex - plantilla maestra UNICA de las guias MILC v3.
%
% La usan las tres familias de guias, cada una con su driver:
%   · Grados 6-11  -> scripts/build-guias-g11.py
%   · Bebras       -> scripts/build-guias-bebras.py
%   · Semillero    -> scripts/build-guias-semillero.py
%
% Antes habia tres copias casi identicas (~400 lineas iguales, 6 distintas):
% cada mejora habia que aplicarla tres veces, y en la practica no se hacia
% (el motor de recursos vivia solo en dos de ellas). Lo que varia entre
% familias esta parametrizado en 6 placeholders que cada driver llena
% (se nombran aqui SIN su sintaxis real: el builder reemplaza por texto plano,
% tambien dentro de los comentarios, y un valor multilinea rompe el preambulo):
%
%   HEADER_IZQ        encabezado izquierdo de pagina
%   HEADER_DER        encabezado derecho de pagina
%   PORTADA_ETIQUETA  rotulo sobre el numero grande (GRADO/MOMENTO/MODULO)
%   PORTADA_SERIE     linea de serie de la portada
%   PORTADA_ID        identificador de la guia en la portada
%   PORTADA_CIRCULO   el circulito de grado (°), o vacio si no aplica
%
% El resto de claves las documenta content/guias/_SCHEMA.md. El script valida
% que no queden placeholders sin reemplazar antes de escribir el archivo final.

\documentclass[11pt,letterpaper]{article}
\usepackage[margin=2.0cm]{geometry}
\usepackage[table]{xcolor}
\usepackage{fontspec}
\usepackage[spanish]{babel}
\usepackage{tikz}
% Librerías TikZ para los diagramas de las guías (bloque `recursos:`):
%  · babel  → babel en español vuelve activos caracteres como «>», y TikZ los usa
%             en sus opciones (>=stealth, ->). Sin esta librería, cualquier
%             diagrama con flechas revienta con «\language@active@arg> has an
%             extra }». Debe ir ANTES de que se dibuje nada.
%  · positioning        → sintaxis «below left=2pt and 3pt».
%  · decorations.pathreplacing → llaves ({brace}) para señalar rangos.
\usetikzlibrary{babel,positioning,decorations.pathreplacing}
\usepackage{graphicx}
% Circuitos: el trabajo de electrónica se hace hoy con micro:bit y sus sensores
% integrados (MakeCode, sin cableado), así que no se necesitan esquemáticos.
% Si algún día entran montajes con protoboard, basta descomentar esta línea:
% los diagramas .tex/.tikz declarados en `recursos:` ya se incrustan solos.
% \usepackage{circuitikz}
\usepackage[most]{tcolorbox}
\usepackage{tabularx}
\usepackage{array}
\usepackage{fancyhdr}
\usepackage{titlesec}
\usepackage[document]{ragged2e}  % bandera a la derecha en todo el cuerpo
\usepackage{enumitem}
\usepackage{microtype}

% Helvetica Neue no trae la flecha U+2192, asi que cada «→» escrito literalmente
% en el YAML se compilaba a NADA, en silencio (462 flechas en 61 guias: rutas del
% tipo «paleta → tipografia → lockup» salian rotas). Mapeamos el caracter a su
% version matematica, que si tiene glifo. Asi el autor puede seguir escribiendo
% «→» comodamente en el YAML.
\usepackage{newunicodechar}
\newunicodechar{→}{\ensuremath{\rightarrow}}
% Mismo problema con los simbolos de marca y vinetas: el «✓» de «una palomita ✓»
% salia como cuadrito vacio en el PDF de todas las guias que lo usaban.
\usepackage{pifont}
\newunicodechar{✓}{\ding{51}}
\newunicodechar{✗}{\ding{55}}
\newunicodechar{✔}{\ding{52}}
\newunicodechar{•}{\textbullet}
\newunicodechar{≥}{\ensuremath{\geq}}
\newunicodechar{≤}{\ensuremath{\leq}}

\setmainfont{Helvetica Neue}
\setsansfont{Helvetica Neue}
\setlength{\parindent}{0pt}
\setlength{\parskip}{6pt}
% Sin particion de palabras: en 6.º-7.º una palabra cortada por guion al final
% de linea («Comuni-cacion») frena la lectura mucho mas que una linea corta.
\hyphenpenalty=10000
\exhyphenpenalty=10000
\pretolerance=2000
\tolerance=3000
\setlength{\emergencystretch}{3em}
\linespread{1.10}
\renewcommand{\arraystretch}{1.25}
\setlist{leftmargin=5mm,itemsep=1mm,topsep=1mm}
\newcolumntype{Y}{>{\RaggedRight\arraybackslash}X}

\definecolor{milcMagenta}{HTML}{E6007E}
\definecolor{milcVino}{HTML}{5A0038}
\definecolor{milcTurquesa}{HTML}{0093A5}
\definecolor{milcVerde}{HTML}{58A923}
\definecolor{milcMostaza}{HTML}{E5B400}
\definecolor{milcOcre}{HTML}{B86600}
\definecolor{milcNegro}{HTML}{241816}
\definecolor{milcGris}{HTML}{F7F7F4}
\definecolor{milcLinea}{HTML}{E8E2DD}
\definecolor{milcGradePrimary}{HTML}{84CC16}
\definecolor{milcGradeDark}{HTML}{4D7C0F}
\definecolor{milcGradeSoft}{HTML}{ECFCCB}
\definecolor{milcGradeAccent}{HTML}{CFE7FF}
\definecolor{milcGradeText}{HTML}{FFFFFF}
\color{milcNegro}

\pagestyle{fancy}
\fancyhf{}
\lhead{\small Tecnología e Informática · Grado 7}
\rhead{\small Guía 23 · MILC v3}
\cfoot{\small\thepage}
\renewcommand{\headrulewidth}{0pt}

\newtcolorbox{titlebox}[1]{
  enhanced,colback=#1,colframe=#1,arc=7mm,boxrule=0pt,
  left=7mm,right=7mm,top=3mm,bottom=3mm,
  before skip=8pt,after skip=8pt,
  fontupper=\sffamily\bfseries\Large\color{white}
}

\newtcolorbox{softbox}[3]{
  enhanced,colback=#2,colframe=#1,borderline west={4pt}{0pt}{#1},
  arc=2mm,boxrule=.4pt,left=4mm,right=4mm,top=3mm,bottom=3mm,
  before skip=6pt,after skip=8pt,title={#3},coltitle=white,
  colbacktitle=#1,fonttitle=\sffamily\bfseries,fontupper=\RaggedRight,
  attach boxed title to top left={xshift=3mm,yshift=-2mm},
  boxed title style={arc=2mm, boxrule=0pt}
}

\newtcolorbox{drawbox}[2]{
  enhanced,colback=white,colframe=#1,arc=4mm,boxrule=1pt,
  left=5mm,right=5mm,top=4mm,bottom=4mm,before skip=8pt,after skip=8pt,
  title={#2},coltitle=white,colbacktitle=#1,fonttitle=\sffamily\bfseries,
  fontupper=\RaggedRight,
  attach boxed title to top left={xshift=4mm,yshift=-2mm},
  boxed title style={arc=3mm, boxrule=0pt}
}

\newtcolorbox{infoband}[1]{
  enhanced,colback=#1!8,colframe=#1!50,arc=2mm,boxrule=.5pt,
  left=4mm,right=4mm,top=2mm,bottom=2mm,before skip=4pt,after skip=4pt,
  fontupper=\small\RaggedRight
}

% Puente narrativo entre secciones (contrato editorial Conéctate).
% Da continuidad: "Acabas de X. Ahora vas a Y porque Z."
% Visualmente: banda izquierda en color del grado, texto en itálica pequeña.
\newtcolorbox{puentebox}{
  enhanced,colback=milcGradeSoft,colframe=milcGradeDark,
  borderline west={3pt}{0pt}{milcGradeDark},
  arc=0mm,boxrule=0pt,left=5mm,right=4mm,top=2.5mm,bottom=2.5mm,
  before skip=4pt,after skip=8pt,
  fontupper=\itshape\small\color{milcNegro}\RaggedRight
}
% La flecha va en modo matematico a proposito: Helvetica Neue NO tiene el glifo
% U+2192, asi que \textrightarrow se compilaba a NADA («Missing character: There
% is no → in font Helvetica Neue») y el puente salia sin flecha. En modo
% matematico la flecha viene de la fuente matematica y si se dibuja.
\newcommand{\puente}[1]{%
  \begin{puentebox}%
    \textcolor{milcGradeDark}{$\rightarrow$}\hspace{2mm}#1%
  \end{puentebox}%
}

\newcommand{\linead}[1]{\noindent\color{milcLinea}\rule{#1}{0.5pt}\color{milcNegro}}
\newcommand{\respuesta}[1]{\par\vspace{2mm}\foreach \n in {1,...,#1}{\noindent\color{milcLinea}\rule{\linewidth}{0.5pt}\par\vspace{5mm}}\color{milcNegro}}
\newcommand{\checkbox}{\textcolor{milcGradeDark}{\fbox{\phantom{x}}}\hspace{5pt}}

\newcommand{\phasecard}[4]{
\begin{tcolorbox}[enhanced,colback=#3,colframe=#2,arc=3mm,boxrule=.5pt,height=3.05cm,valign=center,halign=center,left=1.5mm,right=1.5mm,top=2mm,bottom=2mm]
{\sffamily\bfseries\fontsize{22}{24}\selectfont\textcolor{#2}{#1}}\par
{\sffamily\bfseries\small #4}
\end{tcolorbox}
}

% ── Piezas del rediseno 2026 (legibilidad 6.º-11.º) ───────────────────
% Banda de estacion: reemplaza el titlebox macizo. Titulo a la izquierda,
% dato practico (tiempo/modalidad) a la derecha, en una sola linea.
\newcommand{\milcestacion}[2]{%
\begin{tcolorbox}[enhanced,colback=#1,colframe=#1,arc=2mm,boxrule=0pt,
  left=5mm,right=5mm,top=2.6mm,bottom=2.6mm,before skip=11pt,after skip=7pt]
{\sffamily\bfseries\fontsize{15}{18}\selectfont\color{white}#2}
\end{tcolorbox}}

% Tarjeta de paso numerada: sustituye las tablas «Pilar / Aplicacion», que
% eran formato de consulta docente, no de estudiante. El numero va en un
% circulo sobre el borde para que la secuencia se lea de un vistazo.
\newcommand{\milcpaso}[3]{%
\begin{tcolorbox}[enhanced,colback=white,colframe=milcLinea,arc=1.5mm,boxrule=.9pt,
  left=14mm,right=4mm,top=3mm,bottom=3mm,before skip=4pt,after skip=4pt,
  fontupper=\RaggedRight,
  overlay={\node[anchor=north west,circle,fill=#2,text=white,inner sep=0pt,
    minimum size=7.4mm,font=\sffamily\bfseries\small]
    at ([xshift=3.6mm,yshift=-3.2mm]frame.north west) {#1};}]
#3
\end{tcolorbox}}

% Casilla marcable de tamano util con lapiz (la anterior era \fbox{\phantom{x}},
% de alto variable segun el texto que la seguia).
\newcommand{\milcmarca}{\framebox[3.8mm]{\rule{0pt}{3.4mm}}}

% Fila marcable de lista suelta (checks de la estacion 1).
\newcommand{\milccheckrow}[1]{%
\par\vspace{1.8mm}\noindent
\makebox[6.4mm][l]{\raisebox{-.55mm}{\textcolor{milcNegro}{\milcmarca}}}%
\parbox[t]{\dimexpr\linewidth-6.4mm\relax}{\RaggedRight #1}\par}

% Celda marcable dentro de la rubrica (tabla de autoevaluacion). Misma
% sangria francesa, pero pensada para vivir dentro de una columna X.
\newcommand{\milcopcion}[1]{%
\makebox[5.2mm][l]{\raisebox{-.55mm}{\textcolor{milcNegro}{\milcmarca}}}%
\parbox[t]{\dimexpr\linewidth-5.2mm\relax}{\RaggedRight #1}}

% ── Recursos visuales de la guía ──────────────────────────────────────
% Declarados en el bloque `recursos:` del YAML y emitidos por el builder.
% \guiaFigura   → imagen rasterizada o PDF (foto, carta celeste, montaje).
% \guiaDiagrama → fuente TikZ/circuitikz (.tex) incrustada como vector:
%                 es la vía para circuitos y esquemáticos editables.
% #1 = ruta relativa al .tex · #2 = pie (puede ir vacío)
\newcommand{\guiaFigura}[2]{%
\begin{center}
\includegraphics[width=0.88\linewidth,height=0.38\textheight,keepaspectratio]{#1}\\[1.5mm]
{\footnotesize\itshape\textcolor{milcNegro!75}{#2}}
\end{center}
}

\newcommand{\guiaDiagrama}[2]{%
\begin{center}
\input{#1}\\[1.5mm]
{\footnotesize\itshape\textcolor{milcNegro!75}{#2}}
\end{center}
}

\begin{document}
\thispagestyle{empty}

% ── Portada ───────────────────────────────────────────────────────────
% Antes: un rectangulo de color a pagina completa. Se veia bien en pantalla,
% pero en la fotocopiadora del colegio se comia el toner y salia gris sucio.
% Ahora la banda ocupa el tercio superior y el resto va en blanco.
\begin{tikzpicture}[remember picture,overlay]
  \path[fill=milcGradePrimary]
    (current page.north west) rectangle ([yshift=-10.6cm]current page.north east);

  \node[anchor=north west,text=milcGradeText] at ([xshift=2.0cm,yshift=-1.55cm]current page.north west)
    {\sffamily\bfseries\fontsize{12}{14}\selectfont GRADO};
  \node[anchor=north west,text=milcGradeText,opacity=.85] at ([xshift=2.0cm,yshift=-2.35cm]current page.north west)
    {\sffamily\bfseries\fontsize{8.5}{10}\selectfont SERIE GUÍAS · TECNOLOGÍA E INFORMÁTICA};
  \node[anchor=north east,text=milcGradeText,opacity=.85,align=right] at ([xshift=-2.0cm,yshift=-1.55cm]current page.north east)
    {\sffamily\bfseries\fontsize{9}{12}\selectfont I.E. Sor María Juliana\\Cartago, Valle del Cauca};

  \node[anchor=west,text=milcGradeText] (gradonum) at ([xshift=1.85cm,yshift=-7.1cm]current page.north west)
    {\sffamily\bfseries\fontsize{112}{112}\selectfont 7};
  % El circulito de grado (°) solo aplica a las guias de grado («11°»); en
  % Bebras y Semillero el numero grande es un momento/modulo, no un grado.
  % Se posiciona relativo al nodo `gradonum`, no en coordenadas absolutas.
  \draw[milcGradeText,line width=1.6pt]
    ([xshift=2.0mm,yshift=-4.5mm]gradonum.north east) circle (.15cm);

  \node[anchor=west,text=milcGradeText,text width=11.4cm,align=left] at ([xshift=1.5cm,yshift=0.5cm]gradonum.east)
    {
     {\sffamily\bfseries\fontsize{9}{11}\selectfont\MakeUppercase{Período 3 · Inteligencia Artificial}}\\[3mm]
     {\sffamily\bfseries\fontsize{21}{25}\selectfont\setlength{\baselineskip}{27pt}Tipos de\\[4pt]IA --- estrecha,\\[4pt]general y\\[4pt]modelos\\[4pt]especializados}\\[3.5mm]
     {\sffamily\bfseries\fontsize{15}{18}\selectfont Guía 23-7-TIC}
    };
\end{tikzpicture}

\vspace*{9.4cm}
\vfill

{\sffamily\bfseries\fontsize{8}{10}\selectfont\textcolor{milcGradeDark}{LO QUE TE LLEVAS HOY}}

\begin{tcolorbox}[enhanced,colback=white,colframe=milcNegro,arc=2mm,boxrule=1.6pt,
  left=5mm,right=5mm,top=4mm,bottom=4mm,before skip=3pt,after skip=12pt,
  fontupper=\RaggedRight\fontsize{11}{15.5}\selectfont]
Una \textbf{tabla comparativa de 5 tipos de IA} en el cuaderno con: \textbf{(a) nombre técnico}, \textbf{(b) qué hace}, \textbf{(c) ejemplo concreto (app que ya usas)}, \textbf{(d) ¿es ``estrecha'' o ``general''?}, \textbf{(e) ¿qué NO puede hacer todavía?}. Más \textbf{un análisis personal} de los 3 tipos de IA que más usas tú + reflexión sobre por qué la AGI (Inteligencia Artificial General) todavía no existe.
\end{tcolorbox}

\begin{tcolorbox}[enhanced,colback=milcGris,colframe=milcGris,arc=2mm,boxrule=0pt,
  left=5mm,right=5mm,top=3.5mm,bottom=3.5mm,after skip=12pt]
{\sffamily\bfseries\fontsize{8}{10}\selectfont\textcolor{milcNegro!55}{ESTA GUÍA ES DE}}\\[3mm]
\begin{tabularx}{\linewidth}{X p{3.2cm} p{3.2cm}}
\linead{\linewidth} & \linead{3.0cm} & \linead{3.0cm}\\[-1mm]
{\fontsize{8.5}{10}\selectfont\textcolor{milcNegro!55}{Nombre completo}} &
{\fontsize{8.5}{10}\selectfont\textcolor{milcNegro!55}{Grupo}} &
{\fontsize{8.5}{10}\selectfont\textcolor{milcNegro!55}{Fecha}}\\
\end{tabularx}
\end{tcolorbox}

\vfill

\noindent\textcolor{milcLinea}{\rule{\linewidth}{1pt}}\\[2mm]
{\sffamily\bfseries\fontsize{9.5}{12}\selectfont Autor: Dr. Álvaro Cárdenas Orozco}\\[1mm]
{\sffamily\fontsize{8.5}{11}\selectfont\textcolor{milcNegro!55}{Metodología MILC · apertura ancestral · pensamiento computacional · triángulo Dussel–estoicismo–Floridi}}

\newpage

\section*{Apertura: Tipos de IA --- estrecha, general, modelos especializados}

\begin{softbox}{milcGradeDark}{milcGradeSoft}{Saber ancestral}
\textbf{En el taller de don Lucho el relojero de Cartago, había herramientas especializadas y herramientas generales.} Su \textbf{lupa de aumento} servía solo para una cosa: ver lo pequeño. Su \textbf{cuchillo de mesa} servía para muchas cosas: cortar pan, abrir cajas, raspar. \emph{Don Lucho usaba cada una para lo suyo}: la lupa para ver el engranaje, el cuchillo para tareas variadas. \textbf{Si hubiera tratado de usar la lupa para cortar pan, habría fallado}. La lupa era \emph{especialista} (excelente en una cosa, inútil en todas las demás). El cuchillo era \emph{generalista} (decente en muchas cosas, no excelente en ninguna). \textbf{En IA pasa algo parecido}. Existen IAs especialistas (la mayoría de las que usas hoy) y la idea de IA generalista (que aún no existe). Esa distinción es la clave para entender qué IA es buena para qué.
\end{softbox}

\begin{softbox}{milcTurquesa}{milcGris}{Saber contemporáneo}
En IA, la clasificación más importante es: \textbf{IA estrecha (\emph{narrow AI})} vs \textbf{IA general (\emph{AGI --- Artificial General Intelligence})}. \textbf{IA estrecha:} hace UNA cosa muy bien. ChatGPT genera texto. DALL-E genera imágenes. Google Translate traduce. Spotify recomienda música. \emph{Cada IA es especialista en su tarea}; si la sacas de ahí, falla. \textbf{IA general (AGI):} aún NO existe. Sería una IA que pueda hacer cualquier tarea intelectual humana al mismo nivel que nosotros: aprender un idioma nuevo, resolver problemas de ingeniería, escribir poesía, mantener conversaciones complejas, todo a la vez. \emph{Aún ningún laboratorio del mundo ha logrado AGI}. ChatGPT \emph{parece} general porque conversa de muchos temas, pero técnicamente sigue siendo \emph{especialista en generar texto plausible}. \textbf{Subcategorías de IA estrecha:} (a) modelos de lenguaje (LLM); (b) visión por computador; (c) recomendadores; (d) predicción/clasificación; (e) sistemas expertos. Hoy las conoces.
\end{softbox}

\begin{softbox}{milcMostaza}{milcGris}{Pregunta-puente}
\textit{Si una IA puede ganar al campeón mundial de ajedrez, jugar Go a nivel mundial, recomendar películas y traducir 100 idiomas, ¿\textbf{ya es ``general''}? ¿O es solo muy hábil en cada una de esas cosas por separado?}
\end{softbox}

\begin{softbox}{milcVerde}{milcGris}{Saber y saber hacer}
Al terminar la clase: (1) podrás \textbf{identificar} la diferencia entre IA estrecha y AGI; (2) sabrás \textbf{aplicar} la clasificación a 5 tipos de IA; (3) podrás \textbf{analizar} qué tipos usas más; (4) habrás \textbf{creado} una tabla comparativa de 5 tipos.
\end{softbox}



\small
\begin{tabularx}{\textwidth}{>{\bfseries}p{3.0cm}Y}
\rowcolor{milcGradePrimary!12}Autor & Dr. Álvaro Cárdenas Orozco\\
DBA & Reconoce los fundamentos, usos y límites éticos de la Inteligencia Artificial (MEN, T\&I 7°)\\
Referentes & Saber ancestral del consejero · Modelos de lenguaje · Ética algorítmica y prompting\\
Producto final & Una \textbf{tabla comparativa de 5 tipos de IA} en el cuaderno con: \textbf{(a) nombre técnico}, \textbf{(b) qué hace}, \textbf{(c) ejemplo concreto (app que ya usas)}, \textbf{(d) ¿es ``estrecha'' o ``general''?}, \textbf{(e) ¿qué NO puede hacer todavía?}. Más \textbf{un análisis personal} de los 3 tipos de IA que más usas tú + reflexión sobre por qué la AGI (Inteligencia Artificial General) todavía no existe.\\
\end{tabularx}
\normalsize

\puente{Hoy haces 4 cosas: distingues IA estrecha de AGI, conoces 5 tipos específicos, analizas cuáles usas, armas tabla.}

\milcestacion{milcGradeDark}{Ruta MILC: así vamos a aprender}
\begin{tabularx}{\textwidth}{>{\centering\arraybackslash}X>{\centering\arraybackslash}X>{\centering\arraybackslash}X>{\centering\arraybackslash}X}
\phasecard{1}{milcVerde}{milcGris}{Escucha\\[-1mm]\small Observo, escucho y pregunto.} &
\phasecard{2}{milcTurquesa}{milcGris}{Entender\\[-1mm]\small Sistematizo y ordeno.} &
\phasecard{3}{milcMagenta}{milcGris}{Hacer\\[-1mm]\small Construyo y aplico.} &
\phasecard{4}{milcVino}{milcGris}{Cierre\\[-1mm]\small Evalúo y reflexiono.}
\end{tabularx}


\puente{Empezamos con un test rápido: ¿qué tipo de IA es cada app que conoces?}

\milcestacion{milcVerde}{Estación 1 de 3 · Escucha}

\begin{softbox}{milcVerde}{milcGris}{Escena para escuchar}
\textbf{Actividad 1 · IDENTIFICA --- ¿Qué tipo de IA usa esta app?} (10 min · individual). En el cuaderno responde para cada una de estas 8 apps: \emph{¿qué hace su IA principalmente?}. \emph{(1) TikTok}. \emph{(2) Google Translate}. \emph{(3) ChatGPT}. \emph{(4) Cámara del celular (modo retrato que difumina el fondo)}. \emph{(5) Spotify}. \emph{(6) Banco (cuando rechaza una transacción extraña)}. \emph{(7) Maps (cuando recalcula ruta)}. \emph{(8) Cualquier juego con NPC (personajes no jugables que parecen ``inteligentes'')}. Responde con tu intuición.
\end{softbox}

{\sffamily\bfseries\fontsize{8}{10}\selectfont\textcolor{milcNegro!55}{MARCA CUANDO LO HAYAS HECHO}}
\milccheckrow{Leí las 8 apps con calma.}
\milccheckrow{Pensé qué hace la IA en cada una.}
\milccheckrow{Anoté mi intuición sin buscar.}

\begin{infoband}{milcVerde}


\textbf{Lo que va al cuaderno (Actividad 1 · IDENTIFICA).} Título: «Actividad 1 --- 8 apps y su IA». \textbf{Formato:} lista numerada 1-8 con descripción corta. \textbf{Extensión:} 8 líneas. \textbf{Sabes que terminaste cuando} reconoces patrones (varias apps hacen recomendar, varias generan texto, varias procesan imágenes).
\end{infoband}

\puente{Después aprendes los 5 tipos principales y la distinción estrecha vs general.}

\milcestacion{milcTurquesa}{Estación 2 de 3 · Entender}

\begin{softbox}{milcTurquesa}{milcGris}{Lo que vas a entender}
Las 8 apps caen en \textbf{5 tipos principales de IA estrecha}. Aprende cada uno con sus apps típicas. Y entiende por qué ninguno es \emph{general} (AGI) todavía.
\end{softbox}

\milcpaso{1}{milcTurquesa}{\textbf{Tipo 1 --- Modelos de lenguaje (LLMs).} Generan texto a partir de texto. Conversaciones, resúmenes, traducciones, código. \emph{Apps:} ChatGPT, Claude, Gemini, Copilot. \textbf{Limitación clave:} no \emph{entienden} como humanos --- predicen la palabra más probable basándose en patrones aprendidos de miles de millones de textos. Pueden alucinar (inventar). \textbf{Tipo 2 --- Visión por computador.} Procesa imágenes y videos. Reconoce caras, clasifica objetos, genera imágenes nuevas. \emph{Apps:} desbloqueo facial del celular, modo retrato de cámaras, Google Lens, DALL-E (genera imágenes). \textbf{Limitación clave:} puede confundirse con imágenes adversarias (pequeños cambios invisibles a humanos pueden engañar a la IA).}
\milcpaso{2}{milcTurquesa}{\textbf{Tipo 3 --- Sistemas de recomendación.} Predicen qué te gustaría según tu historial y el de gente similar. \emph{Apps:} TikTok, YouTube, Spotify, Netflix, Amazon. \textbf{Limitación clave:} crean \emph{cámaras de eco} (te muestran solo lo que ya te gusta, no te exponen a cosas nuevas). \textbf{Tipo 4 --- Predicción y clasificación.} Predicen valores numéricos o asignan categorías. \emph{Apps:} Maps (predice tráfico), bancos (predicen fraude), seguros (predicen riesgo de accidente), pronóstico del clima. \textbf{Limitación clave:} solo predicen lo que han visto antes; algo verdaderamente nuevo las confunde.}
\milcpaso{3}{milcTurquesa}{\textbf{Tipo 5 --- Sistemas expertos (Game AI, asistentes específicos).} Tienen reglas y conocimiento específico de un dominio. \emph{Apps:} NPCs de videojuegos (el dragón sabe atacar), asistentes médicos que diagnostican según síntomas, ajedrez profesional, Go (AlphaGo). \textbf{Limitación clave:} fuera del dominio para el que fueron diseñados, no sirven. AlphaGo no puede recomendarte canciones.}
\milcpaso{4}{milcTurquesa}{\textbf{¿Por qué la AGI aún no existe?} La AGI (Artificial General Intelligence) sería una IA que iguale a un humano en CUALQUIER tarea intelectual. \textbf{No existe} en 2026 porque los modelos actuales: (a) Son especialistas en un dominio (ChatGPT en texto, DALL-E en imágenes). (b) No tienen \emph{sentido común} verdadero (cometen errores que un niño no cometería). (c) No tienen \emph{conciencia} ni \emph{intenciones} (procesan, no quieren). (d) No \emph{aprenden} de una sola experiencia como humanos (necesitan miles de ejemplos). \textbf{¿Cuándo llegará AGI? Los expertos están divididos.} Algunos dicen 5-10 años; otros, 50; otros, nunca. \emph{Vivimos en el momento de gran apuesta tecnológica, pero nadie sabe la respuesta segura}.}

\begin{drawbox}{milcTurquesa}{5 tipos + IA estrecha vs general}
\textbf{5 tipos de IA estrecha:} \\[1mm]
\textbf{1.} Modelos de lenguaje (LLMs) --- ChatGPT, Claude. \\[1mm]
\textbf{2.} Visión por computador --- desbloqueo facial, DALL-E. \\[1mm]
\textbf{3.} Recomendadores --- TikTok, YouTube, Spotify. \\[1mm]
\textbf{4.} Predicción/clasificación --- Maps tráfico, bancos antifraude. \\[1mm]
\textbf{5.} Sistemas expertos --- NPCs juegos, asistentes médicos. \\[1mm]
\textbf{IA general (AGI):} aún NO existe.
\end{drawbox}

\begin{softbox}{milcMostaza}{milcGris}{Errores comunes y su corrección}
\textbf{Mitos sobre los tipos de IA.} (1) \emph{``ChatGPT es AGI''} --- Falso. Es especialista en generar texto plausible. (2) \emph{``Las IAs se aburren''} --- Falso. No tienen estados emocionales. (3) \emph{``Si una IA es buena en algo, será buena en todo''} --- Falso. Las IAs estrechas son inútiles fuera de su dominio. (4) \emph{``Los robots de las películas ya existen''} --- Falso. Los robots actuales son muy limitados; los de Hollywood son ciencia ficción. (5) \emph{``La AGI llegará el próximo año seguro''} --- Falso. Nadie sabe cuándo (o si) llegará.
\end{softbox}

\begin{infoband}{milcTurquesa}


\textbf{Lo que va al cuaderno (Actividad 2 · EXPLICA).} Título: «Actividad 2 --- 5 tipos de IA estrecha + AGI». \textbf{Formato:} bloque A con tabla de 5 filas (tipo + ejemplo de app + limitación clave). Bloque B con 3 líneas sobre qué es AGI y por qué aún no existe. \textbf{Extensión:} 10 líneas. \textbf{Sabes que terminaste cuando} podrías clasificar cualquier IA que veas en una de las 5 categorías.
\end{infoband}

\puente{Con todo claro, armas tu tabla comparativa de 5 tipos.}

\milcestacion{milcMagenta}{Estación 3 de 3 · Hacer}

\begin{softbox}{milcMagenta}{milcGris}{Cómo vas a construir tu producto}
\textbf{Actividad 3 · ANALIZA --- Tabla comparativa + análisis personal} (28 min · individual). Vas a armar una tabla completa de los 5 tipos + reflexionar sobre cuáles usas más.
\end{softbox}

\milcpaso{1}{milcMagenta}{\textbf{Paso 1 --- Tabla comparativa de 5 tipos.} En una hoja del cuaderno, tabla de \textbf{5 filas y 5 columnas}. Columnas: \emph{(1) Nombre del tipo}, \emph{(2) Qué hace}, \emph{(3) Ejemplo de app real}, \emph{(4) ¿Estrecha o general?}, \emph{(5) Algo que NO puede hacer}. Llena las 5 filas con los 5 tipos de la sistematización.}
\milcpaso{2}{milcMagenta}{\textbf{Paso 2 --- Análisis: ¿qué tipos uso más?} Debajo de la tabla, subtítulo: \textbf{``Los 3 tipos de IA que más uso''}. Lista numerada de 3 tipos en orden de frecuencia, con justificación corta: \emph{``Uso más [tipo] porque [actividad cotidiana]''}. Ejemplo: \emph{``1. Recomendadores, porque uso TikTok 3 horas al día. 2. Visión por computador, porque desbloqueo el celular 50 veces al día con cara. 3. LLMs, porque uso ChatGPT 5 veces a la semana.''}.}
\milcpaso{3}{milcMagenta}{\textbf{Paso 3 --- Reflexión sobre AGI.} Subtítulo: \textbf{``Por qué la AGI todavía no existe''}. Escribe en 4-5 líneas tu propia explicación: por qué ChatGPT no es AGI, qué le falta a las IAs actuales para ser \emph{generales}, y si crees que llegará la AGI en tu vida (sí, no, no sé).}
\milcpaso{4}{milcMagenta}{\textbf{Paso 4 --- Mapa visual.} En un costado, dibuja un \textbf{círculo grande} con la palabra \emph{IA}. Adentro divide en 5 sectores etiquetados con los 5 tipos (como pastel). Marca con sombreado cuáles usas TÚ regularmente. Esto te da imagen visual de tu uso.}
\milcpaso{5}{milcMagenta}{\textbf{Paso 5 --- Tu predicción.} Subtítulo: \textbf{``Mi predicción para 2030''}. Escribe en 3-4 líneas qué creés que va a pasar con la IA en 4 años: ¿más AGI? ¿más recomendadores? ¿más asistentes? Sé honesto: no copies de internet.}
\milcpaso{6}{milcMagenta}{\textbf{Paso 6 --- Firma y cierre.} Al final firma con nombre y fecha. Tu predicción te servirá: en 4 años revisarás esta hoja y verás si tu intuición fue cierta.}

\begin{drawbox}{milcMagenta}{Antes de cerrar la S3}
\checkbox La tabla tiene 5 tipos con sus 5 columnas. \\[1mm]
\checkbox Identifiqué los 3 tipos que más uso. \\[1mm]
\checkbox Hay reflexión sobre por qué AGI no existe. \\[1mm]
\checkbox Dibujé el mapa visual circular con 5 sectores. \\[1mm]
\checkbox Hice predicción para 2030 con mi voz. \\[1mm]
\checkbox Está firmado con nombre y fecha.
\end{drawbox}

\begin{softbox}{milcMagenta}{milcGris}{Plantilla de trabajo}
\textbf{Estructura.} \textit{Paso 1:} tabla 5x5 (tipo, qué hace, ejemplo, estrecha/general, limitación). \textit{Paso 2:} análisis de los 3 tipos que más uso. \textit{Paso 3:} reflexión sobre AGI. \textit{Paso 4:} mapa visual circular. \textit{Paso 5:} predicción para 2030. \textit{Paso 6:} firma.
\end{softbox}

\begin{infoband}{milcMagenta}


\textbf{Lo que va al cuaderno (Actividad 3 · ANALIZA).} Título: «Actividad 3 --- 5 tipos de IA + análisis personal». \textbf{Formato:} tabla + análisis + reflexión + mapa + predicción. \textbf{Extensión:} 1 página completa. \textbf{Sabes que terminaste cuando} podrías clasificar cualquier IA nueva que aparezca en uno de los 5 tipos.
\end{infoband}

\puente{El producto es esa tabla + análisis personal + reflexión sobre AGI.}

\milcestacion{milcMostaza}{Producto de la sesión}
\textbf{Tabla de 5 tipos de IA + análisis personal + predicción 2030 · firmado}.
\begin{softbox}{milcMostaza}{milcGris}{Criterios de calidad}
(1) La tabla tiene los 5 tipos con las 5 columnas completas. (2) Hay análisis personal de los 3 tipos más usados. (3) Hay reflexión sobre por qué AGI todavía no existe. (4) El mapa visual circular distingue los 5 tipos. (5) Hay predicción personal para 2030.
\end{softbox}

\puente{Antes de salir, verifica que cada tipo tiene su ejemplo de app real.}

\milcestacion{milcVino}{Evaluación liberadora · cinco dimensiones}

\begin{softbox}{milcVino}{milcGris}{Sentido de la evaluación}
Evaluar no es solo revisar una nota. Es reconocer qué aprendí, qué cambió en mí, qué emociones manejé, cómo me relaciono con la ciudad y la comunidad, y qué le dejaré a quien viene detrás.
\end{softbox}

\textbf{Mi reflexión liberadora en cinco dimensiones.} Completa cada frase iniciada:

\small
\begin{tabularx}{\textwidth}{>{\bfseries\RaggedRight\arraybackslash}p{3.4cm}Y>{\RaggedRight\arraybackslash}p{4.6cm}}
\rowcolor{milcVino}\color{white}Dimensión & \color{white}\bfseries Mi respuesta (frame starter) & \color{white}\bfseries Algo que descubrí\\
1. Personal & Lo que más cambió en mí fue \linead{6.5cm} & \linead{4.2cm}\\
2. Emocional & La emoción más fuerte fue \linead{2.5cm} y la regulé \linead{3cm} & \linead{4.2cm}\\
3. Ciudadana & En democracia esto importa porque \linead{6cm} & \linead{4.2cm}\\
4. Local (Cartago) & En mi barrio sería útil para \linead{6.5cm} & \linead{4.2cm}\\
5. Intergeneracional & A mi abuela/abuelo le diría \linead{6.5cm} & \linead{4.2cm}\\
\end{tabularx}
\normalsize

\begin{infoband}{milcVino}
\textbf{Para la próxima sustentación.} Vuelve a esta tabla en seis meses. Verás que las respuestas cambian — no porque lo aprendido cambie, sino porque tú cambias. La evaluación liberadora es una conversación contigo a través del tiempo.
\end{infoband}


\puente{Cerramos con 3 voces sobre por qué entender los tipos te hace mejor usuario.}

\milcestacion{milcGradeDark}{Triángulo de pensamiento · cierre filosófico}

\begin{softbox}{milcVino}{milcGris}{Dussel · lente del nosotros}
\textit{Clasificar lo que te rodea no es burocracia académica: es soberanía intelectual. El que distingue tipos, no se deja engañar.}

{\footnotesize\itshape\color{milcNegro!70}--- formulación de esta guía, al modo de Enrique Dussel. No es una cita textual.}

\textbf{Cuando solo conoces ``la IA'' como categoría única}, todas te parecen iguales. Cuando distingues 5 tipos, ves matices: \emph{TikTok hace recomendaciones, no genera texto; ChatGPT genera texto, no genera imágenes}. Esa distinción es \emph{poder mental}: te coloca en posición de analizar críticamente, no de admirar o temer pasivamente. Don Lucho el relojero distinguía las herramientas de su taller; tú distingues los tipos de IA del mundo digital.

\textbf{¿Antes de hoy, qué tan ``mezcladas'' tenía las IAs en mi cabeza?}
\end{softbox}
\linead{\linewidth}\\[1mm]
\linead{\linewidth}

\begin{softbox}{milcMostaza}{milcGris}{Estoicismo · Marco Aurelio · cuidado interior}
\textit{Lo que la IA no puede hacer hoy, importa tanto como lo que sí puede. Conocer ambos límites es sabiduría.}

{\footnotesize\itshape\color{milcNegro!70}--- formulación de esta guía, al modo de Marco Aurelio. No es una cita textual.}

Hay gente que solo habla de \textbf{lo que la IA puede hacer} (impresionante, ¿no?). Otra gente solo habla de \textbf{lo que aún no puede} (limitada, ¿no?). El sabio conoce \emph{ambos}: lo que puede y lo que no. Eso le permite usarla con criterio --- aprovechar fortalezas, no esperar lo imposible. \emph{Tu tabla con la columna ``limitación clave'' entrena esa sabiduría serena}.

\textbf{¿Tiendo a sobreestimar la IA (esperar que pueda todo) o a subestimarla (no aprovechar lo que puede)?}
\end{softbox}
\linead{\linewidth}\\[1mm]
\linead{\linewidth}

\begin{softbox}{milcTurquesa}{milcGris}{Floridi · lente de la infoesfera}
\textit{Cuando llegue la AGI cambiará el mundo. Pero hoy no existe. Vivir como si ya existiera o como si nunca llegara son dos errores opuestos.}

{\footnotesize\itshape\color{milcNegro!70}--- formulación de esta guía, al modo de Luciano Floridi. No es una cita textual.}

\textbf{Algunos viven como si la AGI ya existiera}: le piden a ChatGPT que decida por ellos. Otros viven como si nunca pudiera existir: rechazan toda IA por miedo. \emph{Los dos son errores opuestos}. \textbf{La actitud sabia: usar las IAs estrechas actuales con criterio, sin esperar magia, y estar atento a cómo evolucionan}. Quizás llegue AGI en 10 años; quizás nunca. Vive con la IA que ES, no con la que imaginas.

\textbf{¿Estoy tratando a ChatGPT como AGI (esperando que sepa todo) o como herramienta estrecha (con criterio sobre lo que sí sabe)?}
\end{softbox}
\linead{\linewidth}\\[1mm]
\linead{\linewidth}

\begin{infoband}{milcNegro}
\textbf{Nota para el docente.} Las tres formulaciones son del autor de la guía, no citas textuales de los pensadores: recogen su lente para pensar el tema, no sus palabras. Si vas a citarlos en clase, remite a la obra original.
\end{infoband}

\milcestacion{milcVino}{Mi autoevaluación · marca una casilla por fila}

{\small
\setlength{\extrarowheight}{2.2pt}
\begin{tabularx}{\textwidth}{>{\RaggedRight\arraybackslash\bfseries}p{3.0cm}YYY}
\rowcolor{milcVino}\color{white}\bfseries Criterio & \color{white}\bfseries Lo logré & \color{white}\bfseries En proceso & \color{white}\bfseries Necesito apoyo\\[.6mm]
Comprensión del tema & \milcopcion{Explico las ideas con mis palabras.} & \milcopcion{Reconozco las ideas, falta precisión.} & \milcopcion{Necesito repasar lo básico.}\\[.8mm]
\rowcolor{milcGris}Pensamiento computacional & \milcopcion{Usé los pilares al armar mi producto.} & \milcopcion{Usé algunos pilares.} & \milcopcion{Necesito releer la estación 2.}\\[.8mm]
Aplicación y producto & \milcopcion{Mi producto cumple los criterios.} & \milcopcion{Está armado, con ajustes pendientes.} & \milcopcion{Aún está en construcción.}\\[.8mm]
\rowcolor{milcGris}Reflexión 5 dimensiones & \milcopcion{Respondí cada una con honestidad.} & \milcopcion{Respondí algunas con detalle.} & \milcopcion{Mis respuestas son muy generales.}\\[.8mm]
Triángulo de pensamiento & \milcopcion{Conecté las 3 voces con mi trabajo.} & \milcopcion{Conecté al menos una.} & \milcopcion{No las relaciono con mi proceso.}\\[.8mm]
\end{tabularx}}

\vspace{3mm}
\textbf{Hoy entendí que:}\\[1mm]
\linead{\linewidth}\\[1mm]
\linead{\linewidth}

\textbf{Mi compromiso para la próxima sesión es:}\\[1mm]
\linead{\linewidth}\\[1mm]
\linead{\linewidth}

\begin{softbox}{milcMostaza}{milcGris}{Cita de despedida · Marco Aurelio}
\textit{``El obstáculo en el camino se vuelve el camino. Lo que estorba al camino se vuelve camino.''}\\[1mm]
Cada error de hoy, bien observado, se convierte en pieza del camino que pisarás mañana.
\end{softbox}

\small
\begin{tabularx}{\textwidth}{>{\bfseries}p{3.6cm}Y}
\rowcolor{milcGradeDark!12}Nombre completo: & \linead{10cm}\\
Fecha: & \linead{4cm} \quad Curso/grupo: \linead{4cm}\\
Mi firma: & \linead{9cm}\\
\end{tabularx}
\normalsize

\begin{infoband}{milcGradeDark}
\textbf{Para la próxima sesión.} Esta semana, cada vez que uses una app con IA, clasifícala mentalmente en uno de los 5 tipos. Te entrenas la mente en distinguir. La próxima clase aprendes \textbf{cómo aprende una IA}: datos, entrenamiento, y el problema clave del \emph{sesgo}.
\end{infoband}

\end{document}
